\scanpage{49}
de sorte que l'extension quadratique est décrite
entièrement par la section $c_0$ de $L^{\otimes2}$. Or comme
\[
\deleted{\unsure{c_0=\mathcal D}}\ \unsure{\mathcal D|\mathscr D|\mathfrak d|\partial}\deleted{\unsure{(|{}_{\xi}}}=b_0^2-4c_0=-4c_0\tag*{\deleted{\unsure{a|$\underline c$}}(91)}
\]
\deleted{\unsure{on a aussi|et aussi}}
la donnée de $c_0$ équivaut à celle de $\mathcal D$. On
trouve ainsi

\uline{Th} Si $2\,1_S$ est inv., la catégorie des Algèbres
quadratiques sur $\mathcal O_S$ (ou \deleted{celle} des rev. quadratiques
\inserted{$\Omega$}\deleted{\unsure{de}} de $S$) équivaut à celle des couples $(L,\mathcal D)$
\uline{où} $L$ est un $\mathcal O_S$-Module inversible, et $\mathcal D$
une section de $L^{\otimes2}$. \deleted{\unsure{Le|Lorsque}} Le fibré affine \deleted{\unsure{de rang 1}}\inserted{enveloppant $\Omega$}
\deleted{\unsure{$E$|$\xi$}} est can\textsuperscript{t} scindé, et son Module des translations
est $L$, de sorte que $\Omega\subset\check{\mathbb V}(L)$\deleted{\unsure{)|$\xi$}}. Ceci
dit, si \deleted{\unsure{$\Omega$ s'identifie|l'Algèbre}}\inserted{$\Omega(\mathcal D)$ est le} \deleted{\unsure{le}} sous-schéma de $\check{\mathbb V}(L)$
des sections $x$ satisfaisant $x^2=\mathcal D$, on a \deleted{\unsure{$\Omega=\frac12$}}
\[
\Omega=\frac12\Omega(\mathcal D)\tag*{(92)}
\]

Il reste à vérifier cette dernière assertion.
Or pour un scindage $s$ quelconque, si $\underline b,\underline c$ sont
les sections associées de $L,L^{\otimes2}$, On a vu
dans \S\,2 (variante marginale) que $\Omega$ s'identifie
au sous-schéma de $\check{\mathbb V}(L)$ d'équation
\[
x^2-bx+c=0\tag*{\unsure{(6)|(5)|(*)}}
\]
Ce qui ici, \deleted{\unsure{si}} comme $b=b_0=0$, s'écrit
\[
x^2=\frac14\mathcal D
\]
ou encore
\deleted{\illegible{1}}
\[
(2x)^2=\mathcal D
\]
% archive sheet number: 48
